% Part: incompleteness
% Chapter: arithmetization-syntax
% Section: introduction

\documentclass[../../../include/open-logic-section]{subfiles}

\begin{document}

\olfileid{inc}{art}{int}
\olsection{પરિચય}

સંગણનીયતા અને તર્કને જોડવા માટે તર્કના
પદાર્થો—ચિહ્નો, પદો, !!{formula}s અને
!!{derivation}s—તેમના પરની ક્રિયાઓ તથા
તેમના ગુણધર્મો અને સંબંધો વિશે એવી રીતે
વાત કરવી જોઈએ કે તેમને સંગણકીય રીતે
સંભાળી શકાય. ચિહ્નો, ચિહ્નોના અનુક્રમો
અને તેમાંથી બનેલા બીજા પદાર્થો પરનાં
સંગણનીય વિધેયો અને સંબંધો સીધાં લઈને
આવું કરી શકાય. તાર્કિક વાક્યરચનાના
બધા પદાર્થો સાન્ત છે અને ચિહ્નોના
!!a{enumerable} ગણમાંથી બને છે; તેથી
સંગણનાના કેટલાક નિદર્શોમાં આ શક્ય છે.
પરંતુ બીજા નિદર્શો—જેમ કે પુનરાવર્તી
વિધેયો—સંખ્યાઓ અને તેમના સંબંધો તથા
વિધેયો સુધી મર્યાદિત છે. વધુમાં, અંતે
આપણે અમુક સિદ્ધાંતોની અંદર પણ વાક્યરચના
સંભાળવી છે, ખાસ કરીને અંકગણિતની ભાષામાં
રચાયેલા સિદ્ધાંતોની અંદર. આવી સ્થિતિમાં
વાક્યરચનાને \emph{અંકગણિતીકૃત} કરવી
જરૂરી છે: એટલે વાક્યરચનાના પદાર્થો,
તેમના પરની ક્રિયાઓ અને તેમના સંબંધોને
અનુક્રમે સંખ્યાઓ, અંકગણિતીય વિધેયો અને
અંકગણિતીય સંબંધો તરીકે નિરૂપવા. આ
વિચાર લાઇબનિટ્ઝ સુધી પહોંચે છે: વાક્યરચનાના
પદાર્થોને સંખ્યાઓ આપવી.

ચિહ્નોને તેમના ``સંકેત નંબર'' તરીકે સંખ્યાઓ
આપવી પ્રમાણમાં સરળ છે. કેટલાક ચિહ્નો
થોડો પડકાર ઊભો કરે છે: દાખલા તરીકે,
અનંત !!{variable}s છે અને દરેક
સ્થાનિકતા~$n$ માટે પણ અનંત
!!{function}s છે. છતાં ચિહ્નોને પદ્ધતિસર
અલગ સંખ્યાઓ આપી શકાય; જેમ કે
$\Obj v_2$ અને~$\Obj v_3$ના સંકેત
નંબર જુદા હોય. ચિહ્નોના અનુક્રમો—પદો
અને !!{formula}s જેવા—વધુ મોટો પડકાર
છે. પરંતુ સંખ્યાઓના અનુક્રમોને શુદ્ધ
અંકગણિતીય રીતે સંભાળી શકીએ (દાખલા તરીકે,
અવિભાજ્ય સંખ્યાઓની ઘાતો વડે અનુક્રમોને
સંકેતબદ્ધ કરીને), તો એક-એક ચિહ્નના
સંકેતબદ્ધીકરણને ચિહ્નોના અનુક્રમો સુધી,
અને પછી !!{formula}sના અનુક્રમો અથવા
!!{derivation}s જેવી બીજી ગોઠવણો સુધી
વિસ્તારી શકીએ. આ વિસ્તૃત સંકેતબદ્ધીકરણને
``ગડેલ-સંખ્યાકરણ'' કહે છે. દરેક પદ,
!!{formula} અને !!{derivation}ને એક
ગડેલ-સંખ્યા અપાય છે.

ચિહ્નોના અનુક્રમોને તેમના સંકેત નંબરોના
અનુક્રમો તરીકે સંકેતબદ્ધ કરીને અને આવા
અનુક્રમોનું સંકેતબદ્ધીકરણ સંગણનીય વિધેયોથી
સંભાળી શકાય એવું પસંદ કરીને, ગડેલ-સંખ્યાઓ
પર પણ સંગણનીય વિધેયો વાપરી શકીએ છીએ.
વ્યવહારમાં જરૂરી બધા વિધેયો આદિમ
પુનરાવર્તી હશે. દાખલા તરીકે, અનુક્રમની
લંબાઈ અને તેના સંકેત નંબર પરથી તેનો
$i$મો ઘટક ગણવાના બંને વિધેયો આદિમ
પુનરાવર્તી છે. અનુક્રમનો સંકેત નંબર,
દાખલા તરીકે, !!a{formula}~$!A$ની
ગડેલ-સંખ્યા હોય, તો !!a{formula}ની
લંબાઈ અને તે !!a{formula}માંના $i$મા ચિહ્નનો સંકેત
નંબર પણ~$!A$ની ગડેલ-સંખ્યામાંથી
ગણી શકાય છે. યોગ્ય રીતે રચાયેલા પદની
અથવા યોગ્ય !!{derivation}ની ગડેલ-સંખ્યા
હોવાનો ગુણધર્મ આદિમ પુનરાવર્તી છે તે
સાબિત કરવું થોડું કઠિન છે. છતાં આવું
શક્ય છે, કારણ કે આપણે જોતા અનુક્રમો—પદો,
!!{formula}s અને !!{derivation}s—આગમનાત્મક
રીતે વ્યાખ્યાયિત છે.

ઉદાહરણ તરીકે, પ્રતિસ્થાપનની ક્રિયા લો.
$!A$ સૂત્ર, $x$ ચલ અને $t$ પદ હોય,
તો~$\Subst{!A}{t}{x}$ એ~$!A$માં
$x$ની દરેક મુક્ત ઘટનાને~$t$થી
બદલવાનું પરિણામ છે. હવે ધારો કે~$!A$,
$x$, $t$ને અનુક્રમે ગડેલ-સંખ્યાઓ
$k$, $l$ અને~$m$ આપી છે. એ જ
યોજના~$\Subst{!A}{t}{x}$ને પણ
ગડેલ-સંખ્યા આપે છે, માનો~$n$.
$k$, $l$ અને~$m$માંથી~$n$માં
જતો આ નકશો પ્રતિસ્થાપનની ક્રિયાનો
અંકગણિતીય સમકક્ષ છે. પ્રતિસ્થાપનની
ક્રિયા $!A$, $x$, $t$ને
$\Subst{!A}{t}{x}$માં મોકલે છે;
અંકગણિતીકૃત પ્રતિસ્થાપન વિધેય
ગડેલ-સંખ્યાઓ $k$, $l$, $m$ને
ગડેલ-સંખ્યા~$n$માં મોકલે છે.
આ વિધેય આદિમ પુનરાવર્તી છે તે આપણે
જોઈશું.

વાક્યરચનાનું અંકગણિતીકરણ માત્ર અમૂર્ત
રસની વાત નથી. અલબત્ત, અંકગણિતની ભાષા
જેવી ભાષા પાસે ભાષાઓ વિશે ``વાત કરવાની''
અલગ વ્યવસ્થા ન હોવા છતાં અભિવ્યક્તિઓના
જટિલ ગુણધર્મોને ઔપચારિક બનાવી શકે છે,
એ મૂળે એક નોંધપાત્ર સમજ હતી. પછી એ
ભાષામાં રચાયેલો પિયાનો અંકગણિત જેવો
સિદ્ધાંત પોતાની ભાષા વિશે શું \emph{સાબિત}
કરી શકે—જેમ કે !!{sentence}s સાબિત
થઈ શકે કે સત્ય છે—એ પૂછવું સહેલું બને.
આથી ગડેલનાં (અસાબિતી વિશેનાં) અને
ટાર્સ્કીનાં (સત્યની અવ્યાખ્યેયતા વિશેનાં)
પ્રસિદ્ધ મર્યાદા-પ્રમેયો તરફ જઈએ છીએ.
પરંતુ વાક્યરચનાને અંકગણિતીકૃત કરવાની
યુક્તિ સંગણનીયતા-સિદ્ધાંતનાં મહત્વનાં
પરિણામો સાબિત કરવા માટે પણ ઉપયોગી છે;
દાખલા તરીકે, સિદ્ધાંતોની સંગણકીય શક્તિ
અથવા સંગણનીયતાનાં જુદાં નિદર્શો વચ્ચેનો
સંબંધ. વાક્યરચનાનું અંકગણિતીકરણ બીજા
પદાર્થો અને ગુણધર્મોને અંકગણિતીકૃત કરવાનો
નમૂનો આપે છે. દાખલા તરીકે, સંરૂપણો
અને સંગણનાઓ—માનો ટ્યુરિંગ મશીનની—પણ
એ રીતે અંકગણિતીકૃત કરી શકાય. તેથી
એક નિદર્શની સંગણનાને, જેમ કે ટ્યુરિંગ
મશીનની, બીજા નિદર્શમાં, જેમ કે પુનરાવર્તી
વિધેયોમાં, અનુકરી શકાય છે.

\end{document}
